\begin{definition}[The fractional energy form]
For $v, w \in L^2(\gamma_n)$ with Hermite coefficients $v_\alpha, w_\alpha$
relative to \noderef{HermiteBasis}, define the real number
\[
  \mathcal{E}(v,w) = \sum_{\alpha \neq 0} |\alpha|^{1/2}\, v_\alpha\, w_\alpha
\]
whenever the family on the right is summable, and $\mathcal{E}(v,w) = 0$
otherwise (a convention never invoked in any statement below).  The family is
absolutely summable as soon as both $v$ and $w$ have finite energy in the sense of
\noderef{Energy}: by Cauchy--Schwarz,
\[
  \sum_{\alpha \neq 0} |\alpha|^{1/2}\,|v_\alpha w_\alpha|
    \leq \Bigl(\sum_{\alpha\neq0}|\alpha|^{1/2}|v_\alpha|^2\Bigr)^{1/2}
         \Bigl(\sum_{\alpha\neq0}|\alpha|^{1/2}|w_\alpha|^2\Bigr)^{1/2}
    = \mathcal{E}(v)^{1/2}\,\mathcal{E}(w)^{1/2}.
\]
On $\mathcal{V}$ the map $(v,w) \mapsto \mathcal{E}(v,w)$ is therefore a
well-defined symmetric bilinear form with $\mathcal{E}(v,v) = \mathcal{E}(v)$,
and it satisfies the Cauchy--Schwarz bound just displayed.  It is the energy form
of the reference paper, equal to $\langle L^{1/4}v, L^{1/4}w\rangle$ on
$\mathcal{V} = \operatorname{Dom}(L^{1/4})$.  It vanishes whenever either argument
is constant, since the index $\alpha = 0$ is omitted.
\end{definition}
